\section{Background and Related Work}\label{sec:related}

\subsection{Datasets}\label{sec:data}
\emph{MOT17} \cite{milan2016mot16} contains people filmed with static and moving cameras; as is
common, the second half of its seven training sequences, val-half (2,652 frames), is used for
evaluation \cite{zhang2022bytetrack}, with the published YOLOX-X detections \cite{ge2021yolox} at the
two published emission floors (lowest emitted scores) of 0.01 and 0.1. To study score changes
directly, these scores were also changed before the tracker by three order-preserving transforms:
cubing ($s^3$), halving ($0.5\,s$), and logit temperature $T=2$, which divides the logit
of each score (Sect.~\ref{sec:problem}) by $T$. \emph{KITTI tracking} \cite{geiger2012kitti} has 21
training sequences (8,008 frames) recorded from a car, scored on cars and pedestrians, and is used
with two other detectors, YOLOv8n \cite{jocher2023yolov8} and RT-DETR-L \cite{zhao2024rtdetr}.
\emph{UAVDT} \cite{du2018uavdt} contains drone videos of vehicles; its 20 test sequences (16,592
frames) are used with ByteTrack, YOLOv8n, and RT-DETR-L. MOT17 with four trackers, KITTI, and the
score transforms were development data. Two predeclared external trackers, PD-SORT and Hybrid-SORT,
were run once with the frozen configuration on their authors' MOT17 val-half detections, and the
held-out UAVDT test split was run once after a written plan. Each external host was checked against
its authors' number: EXACT within 0.2 HOTA points, CLOSE within 1.0.

\subsection{Evaluation metrics}\label{sec:metrics}
HOTA \cite{luiten2021hota}, the main metric, balances finding the objects and keeping their
identities,
\begin{equation}
\mathrm{HOTA}=\frac{1}{19}\sum_{\alpha}\sqrt{\mathrm{DetA}_\alpha\,\mathrm{AssA}_\alpha},
\label{eq:hota}
\end{equation}
where $\mathrm{DetA}_\alpha$ and $\mathrm{AssA}_\alpha$ are the detection and association accuracy
at the localization threshold $\alpha=0.05,0.10,\dots,0.95$.

MOTA \cite{bernardin2008clear} and IDF1 \cite{ristani2016idf1} are
\begin{align}
\mathrm{MOTA}&=1-\frac{\mathrm{FN}+\mathrm{FP}+\mathrm{IDS}}{\mathrm{GT}},\label{eq:mota}\\
\mathrm{IDF1}&=\frac{2\,\mathrm{IDTP}}{2\,\mathrm{IDTP}+\mathrm{IDFP}+\mathrm{IDFN}},\label{eq:idf1}
\end{align}
where FN are missed objects (false negatives), FP false positives, IDS identity switches, and GT the
number of ground-truth boxes; IDTP, IDFP, and IDFN count boxes after each ground-truth track is
matched to at most one output track.

All metrics are computed with TrackEval (commit 12c8791) and the official ground truth
\cite{dendorfer2021motchallenge}; KITTI HOTA is averaged over cars and pedestrians, and on UAVDT all
vehicles form one class and predictions in ignored regions are removed. HOTA, MOTA, and IDF1 are in
percent and higher is better ($\uparrow$); IDS is a count and lower is better ($\downarrow$);
differences are in points. Each tracker runs alone (Host) and with \scl{} (Host+\scl{}) on the same
sequences and detections, and Change is Host+\scl{} minus Host. Paired bootstrap resampling of whole
sequences \cite{efron2016casi}, 10,000 times with seed 42, gives a 95\% interval for each
change, which is significant (bold in the tables) when its interval does not include zero.

\subsection{Related work}\label{sec:survey}
\emph{Tracking by detection.}\label{sec:hosts} SORT matches Kalman-filter predictions to new
detections by their overlap \cite{bewley2016sort}; ByteTrack added a second step for low-score
detections \cite{zhang2022bytetrack}; OC-SORT uses one score threshold \cite{cao2023ocsort}. Their
thresholds were tuned for one detector, usually YOLOX \cite{ge2021yolox}. Many recent trackers improve
association with trained or domain-specific parts
\cite{li2024fdbtrack,bao2024graphmatching,zhang2025sinkhorn,ding2025remot,liang2024motiontrack,sun2026bsr},
which change or retrain the tracker; \scl{} keeps both unchanged. Seven published trackers serve as
hosts with the authors' code and configuration: ByteTrack, BoT-SORT \cite{aharon2022botsort},
OC-SORT, SparseTrack \cite{liu2025sparsetrack}, BoostTrack \cite{stanojevic2024boosttrack}, PD-SORT
\cite{wang2025pdsort}, and Hybrid-SORT \cite{yang2024hybridsort}. ByteTrack is used with its authors'
MOT17 thresholds (official setting) and with the Ultralytics defaults, which also serve BoT-SORT
(library setting).

\emph{Adaptive thresholds.} Adaptive ByteTrack sets ByteTrack's threshold from each frame
\cite{ma2023adaptivebytetrack}, AdapTrack adapts its matching thresholds \cite{shim2024adaptrack},
and ACT sets the high and low thresholds from the first quartile and the number of the current
scores \cite{act2026adaptive}. ConfTrack weights the Kalman update and the matching cost by the
score \cite{jung2024conftrack}, and PlugTrack blends motion predictors with a trained network
\cite{kim2026plugtrack}.

\emph{Comparison.} Table~\ref{tab:compare} compares \scl{} with these methods. The
adaptive-thresholding methods act in every frame, read only the current frame or each box, and are
built into one tracker; calibration needs labelled data. \scl{} first decides from past scores
whether to act at all, serves unchanged single-stage and two-stage trackers through their published
settings, and needs neither labels nor trained parts (Sect.~\ref{sec:choice}).

\begin{table*}[!tp]
\caption{\scl{} compared with adaptive-thresholding and score-calibration methods: what each method
changes, the information it uses, whether it needs labels or trained parts, when it acts, and where
it works. ReID: re-identification network.}
\label{tab:compare}
\centering\footnotesize
\setlength{\tabcolsep}{2.5pt}
\begin{tabular}{@{}>{\raggedright\arraybackslash}p{26mm}>{\raggedright\arraybackslash}p{33mm}>{\raggedright\arraybackslash}p{32mm}>{\raggedright\arraybackslash}p{17mm}>{\raggedright\arraybackslash}p{19mm}>{\raggedright\arraybackslash}p{24mm}@{}}
\toprule
Method & What is adapted & Information used & Labels or trained parts & When it acts & Where it works\\
\midrule
\multicolumn{6}{@{}l}{\textit{Adaptive thresholding}}\\
Adaptive ByteTrack \cite{ma2023adaptivebytetrack} & high/low score split & scores of the current frame & none & every frame & inside ByteTrack\\
AdapTrack \cite{shim2024adaptrack} & matching thresholds & current frame & ReID network & every frame & its own tracker\\
ACT \cite{act2026adaptive} & high and low thresholds, from preset levels & first quartile and number of the current scores & none & every frame & inside the tracker\\
ConfTrack \cite{jung2024conftrack} & Kalman update, matching cost & score of each box & ReID network & every box & its own tracker\\
\addlinespace[2pt]
\multicolumn{6}{@{}l}{\textit{Score calibration}}\\
K{\"u}ppers et al.\ \cite{kuppers2020detcalib} & scores mapped to probabilities & labelled data & calibration map & every score & detector output\\
\addlinespace[2pt]
\multicolumn{6}{@{}l}{\textit{This paper}}\\
\scl{} & candidates, scores, thresholds, matching tolerance & logits of past frames; host contract & none & noisy streams (clean: capped thresholds) & between a detector and unchanged trackers\\
\bottomrule
\end{tabular}
\end{table*}



\emph{Score calibration and adaptive inference.} Detector families such as YOLOX, YOLOv8, and
RT-DETR give different score distributions, and calibration methods learn a mapping from labelled
data \cite{kuppers2020detcalib}. Otsu's method picks, without labels, the threshold that best splits a
histogram into two groups \cite{goh2018otsu}; \scl{} applies it to the logits of recent scores. Dynamic
networks learn their settings for each model \cite{han2022dynamic}, and \priorctl{}
sets the detector operating point from image and track cues. \scl{} instead
reads the scores and acts on the tracker's input and thresholds.
